\setcounter{section}{14}






  \zwischenueberschrift{Algebraische Funktionen auf Varietäten}

Was ist ein Morphismus zwischen zwei affin-algebraischen Mengen $V$ und $W$? Wir betrachten zuerst die Situation, wo

\mavergleichskette
{\vergleichskette
{ W
}
{ = }{  {\mathbb A}^{1}_{K}
}
{  }{
}
{    }{
}
{   }{
}
}
{}{}{}
die affine Gerade ist. Sei

\mavergleichskette
{\vergleichskette
{ V
}
{ = }{ V(  {\mathfrak a} )
}
{ \subseteq }{ { {\mathbb A}_{ K }^{ n  } }
}
{    }{
}
{   }{
}
}
{}{}{}
als abgeschlossene Teilmenge eines affinen Raumes gegeben. Dann liefert jedes Polynom

\mavergleichskette
{\vergleichskette
{ F
}
{ \in }{  K[X_1  , \ldots ,  X_n]
}
{  }{
}
{    }{
}
{   }{
}
}
{}{}{}
eine Abbildung
\maabb {F} { { {\mathbb A}_{ K }^{ n  } } } { {\mathbb A}^{1}_{K}  = K
} {}
und damit durch Einschränkung auch eine Abbildung auf $V$. Das haben wir schon bei der Definition des Koordinatenrings betrachtet. Ebenso liefert ein Element

\mavergleichskette
{\vergleichskette
{ F
}
{ \in }{ R
}
{  }{
}
{    }{
}
{   }{
}
}
{}{}{}
in einer endlich erzeugten $K$-Algebra $R$ eine Funktion auf $K\!-\!\operatorname{Spek}\, \left(  R  \right)$, nämlich
\maabbeledisp {} { K\!-\!\operatorname{Spek}\, \left(  R  \right) } { {\mathbb A}^{1}_{K}
} { P } { F(P)
} {.}
Dies ist auch die
\definitionsverweis {Spektrumsabbildung}{}{,}
die nach

Proposition 12.8  (2)

zum
\definitionsverweis {Einsetzungshomomorphismus}{}{}
\maabbeledisp {} {K[T]  } { R
} {T} { F
} {,}
gehört.

Für die offenen Mengen

\mavergleichskette
{\vergleichskette
{  D(F)
}
{ \cong }{ K\!-\!\operatorname{Spek}\, \left(  R_F  \right)
}
{  }{
}
{    }{
}
{   }{
}
}
{}{}{}
ist $1/F$
nach
Satz 13.4

eine wohldefinierte Funktion. Wir werden allgemein für eine Zariski-offene Menge

\mavergleichskette
{\vergleichskette
{U
}
{ \subseteq }{V
}
{  }{
}
{    }{
}
{   }{
}
}
{}{}{}
erklären, was eine algebraische Funktion auf $U$ ist. Die folgende Definition ist so strukturiert, dass die Bedingung \anfuehrung{algebraisch}{} eine \stichwort {lokale Eigenschaft} {} ist.


\inputdefinition
{}
{





Es sei $K$ ein
\definitionsverweis {algebraisch abgeschlossener Körper}{}{,}
$R$ eine
$K$-\definitionsverweis {Algebra von endlichem Typ}{}{}
und sei

\mavergleichskette
{\vergleichskette
{ V
}
{ = }{ K\!-\!\operatorname{Spek}\, \left(  R  \right)
}
{  }{
}
{    }{
}
{   }{
}
}
{}{}{}
das
$K$-\definitionsverweis {Spektrum}{}{}
von $R$. Es sei

\mavergleichskette
{\vergleichskette
{ P
}
{ \in }{ V
}
{  }{
}
{    }{
}
{   }{
}
}
{}{}{}
ein Punkt,

\mavergleichskette
{\vergleichskette
{ U
}
{ \subseteq }{ V
}
{  }{
}
{    }{
}
{   }{
}
}
{}{}{}
eine
\definitionsverweis {Zariski-offene Menge}{}{}
mit

\mavergleichskette
{\vergleichskette
{ P
}
{ \in }{ U
}
{  }{
}
{    }{
}
{   }{
}
}
{}{}{}
und es sei
\maabb {f} { U } { {\mathbb A}^{1}_{ K } = K
} {}
eine Funktion. Dann heißt $f$ \definitionswort {algebraisch}{}
\zusatzklammer {oder \definitionswort {regulär}{} oder \definitionswort {polynomial}{}} {} {}
im Punkt $P$, wenn es Elemente

\mavergleichskette
{\vergleichskette
{ G,H
}
{ \in }{ R
}
{  }{
}
{    }{
}
{   }{
}
}
{}{}{}
gibt mit

\mavergleichskette
{\vergleichskette
{ P
}
{ \in }{ D(H)
}
{ \subseteq }{ U
}
{    }{
}
{   }{
}
}
{}{}{}
und mit

\mathdisp {f(Q)= { \frac{  G(Q) }{ H(Q) } } \text{ für alle } Q \in D(H)} { . }

Die Funktion $f$ heißt \definitionswort {algebraisch}{}
\zusatzklammer {oder \definitionswort {algebraisch auf $U$ }{}} {} {,}
wenn $f$ in jedem Punkt von $U$ algebraisch ist.





}

Natürlich definiert jedes Element

\mavergleichskette
{\vergleichskette
{ f
}
{ \in }{ R
}
{  }{
}
{    }{
}
{   }{
}
}
{}{}{}
eine algebraische Funktion auf jeder offenen Teilmenge des $K$-Spektrums. Es ist aber im Allgemeinen eher schwierig, die algebraischen Funktionen übersichtlich zu beschreiben.

\inputbemerkung
{}
{





In der

Definition 14.1

ist die Bedingung

\mavergleichskette
{\vergleichskette
{ D(H)
}
{ \subseteq }{ U
}
{  }{
}
{    }{
}
{   }{
}
}
{}{}{}
nicht wichtig. Wenn es eine Beschreibung für $f$ mit

\mavergleichskette
{\vergleichskette
{ f
}
{ = }{ G/H
}
{  }{
}
{    }{
}
{   }{
}
}
{}{}{}
auf
\mathl{D(H)}{} mit

\mavergleichskette
{\vergleichskette
{ P
}
{ \in }{ D(H)
}
{  }{
}
{    }{
}
{   }{
}
}
{}{}{}
gibt, so betrachtet man ein $H'$ mit

\mavergleichskette
{\vergleichskette
{ P
}
{ \in }{ D(H')
}
{  }{
}
{    }{
}
{   }{
}
}
{}{}{,}

\mavergleichskette
{\vergleichskette
{ D(H')
}
{ \subseteq }{ U
}
{  }{
}
{    }{
}
{   }{
}
}
{}{}{.}
Dann kann man zu

\mavergleichskette
{\vergleichskette
{ D(H) \cap D(H')
}
{ = }{ D( H H' )
}
{  }{
}
{    }{
}
{   }{
}
}
{}{}{}
übergehen, und dort die Darstellung

\mavergleichskette
{\vergleichskette
{ f
}
{ = }{ (GH')/(HH')
}
{  }{
}
{    }{
}
{   }{
}
}
{}{}{}
betrachten.

Wenn es im Punkt $P$ eine Bruchdarstellung für $f$ als

\mavergleichskette
{\vergleichskette
{ f
}
{ = }{ G/H
}
{  }{
}
{    }{
}
{   }{
}
}
{}{}{}
gibt, so kann man diese Darstellung für alle Punkte aus
\mathl{D(H)}{} verwenden. D.h. $f$ ist auf der ganzen offenen Menge
\mathl{D(H)}{} algebraisch. Insbesondere muss man nicht mit unendlich vielen verschiedenen Darstellungen arbeiten, sondern man kann sich auf die
\zusatzklammer {endlich vielen} {} {}
Darstellungen
\mathl{G_i/H_i}{} zu einer Überdeckung

\mavergleichskette
{\vergleichskette
{ U
}
{ = }{ \bigcup_{i \in I} D(H_i)
}
{  }{
}
{    }{
}
{   }{
}
}
{}{}{}
beschränken.

Bei

\mavergleichskette
{\vergleichskette
{ K
}
{ = }{ {\mathbb C}
}
{  }{
}
{    }{
}
{   }{
}
}
{}{}{}
ist eine algebraische Funktion auch stetig bezüglich der metrischen Topologie, und bei

\mavergleichskette
{\vergleichskette
{ R
}
{ = }{ {\mathbb C}[X_1 , \ldots , X_n]
}
{  }{
}
{    }{
}
{   }{
}
}
{}{}{}
ist sie \stichwort {holomorph} {.}





}

\inputbeispiel{}
{





Es sei

\mavergleichskette
{\vergleichskette
{ V
}
{ = }{ V(WX-ZY)
}
{ \subseteq }{ { {\mathbb A}_{  K  }^{  4   } }
}
{    }{
}
{   }{
}
}
{}{}{}
und sei

\mavergleichskette
{\vergleichskette
{ U
}
{ = }{ D(X,Y)
}
{ = }{ D(X) \cup D(Y)
}
{ \subset   }{ V
}
{   }{
}
}
{}{}{}
die durch $X$ und $Y$ definierte Zariski-offene Menge. Auf $U$ ist die durch

\mavergleichskettedisp
{\vergleichskette
{ f
}
{ =} { { \frac{   Z  }{  X } }
}
{ =} { { \frac{   W  }{  Y } }
}
{ } {
}
{ } {
}
}
{}{}{}
definierte Funktion algebraisch. Die beiden rationalen Darstellungen liefern offenbar eine algebraische Funktion auf den beiden offenen Teilmengen $D(X)$ und $D(Y)$. Damit es eine Funktion auf $U$ definiert, muss sichergestellt werden, dass die Brüche auf dem Durchschnitt, also auf

\mavergleichskette
{\vergleichskette
{ D(X) \cap D(Y)
}
{ = }{ D(XY)
}
{  }{
}
{    }{
}
{   }{
}
}
{}{}{,}
die gleichen Funktionswerte haben. Es sei also

\mavergleichskette
{\vergleichskette
{ Q
}
{ = }{ (w,x,y,z)
}
{ \in }{ D(XY)
}
{    }{
}
{   }{
}
}
{}{}{,}

\mavergleichskette
{\vergleichskette
{ Q
}
{ \in }{ V
}
{  }{
}
{    }{
}
{   }{
}
}
{}{}{.}
D.h.

\mavergleichskette
{\vergleichskette
{ x,y
}
{ \neq }{ 0
}
{  }{
}
{    }{
}
{   }{
}
}
{}{}{}
und

\mavergleichskette
{\vergleichskette
{ wx
}
{ = }{ zy
}
{  }{
}
{    }{
}
{   }{
}
}
{}{}{.}
Dann ist aber sofort

\mavergleichskettedisp
{\vergleichskette
{ \frac{Z}{X}(Q)
}
{ =} { \frac{z}{x}
}
{ =} { \frac{w}{y}
}
{ =} { \frac{W}{Y}(Q)
}
{ } {
}
}
{}{}{.}





}





\inputfaktbeweis
{K-Spektrum/Algebraisch abgeschlossener Körper/Algebraische Funktion auf offener Menge/Ring/Fakt}
{Lemma}
{}
{



\faktsituation {}
\faktvoraussetzung {Es sei $K$ ein
\definitionsverweis {algebraisch abgeschlossener Körper}{}{,}
$R$ eine
$K$-\definitionsverweis {Algebra von endlichem Typ}{}{}
und sei

\mavergleichskette
{\vergleichskette
{V
}
{ = }{ K\!-\!\operatorname{Spek}\, \left(  R  \right)
}
{  }{
}
{    }{
}
{   }{
}
}
{}{}{}
das
$K$-\definitionsverweis {Spektrum}{}{}
von $R$. Es sei

\mavergleichskette
{\vergleichskette
{ U
}
{ \subseteq }{ V
}
{  }{
}
{    }{
}
{   }{
}
}
{}{}{}
eine
\definitionsverweis {Zariski-offene}{}{}
Menge.}
\faktfolgerung {Dann bildet die Menge der
\definitionsverweis {algebraischen Funktionen}{}{}
auf $U$ einen Unterring
\zusatzklammer {und zwar eine $K$-Unteralgebra} {} {}
des Rings der Funktionen von $U$ nach $K$}
\faktzusatz {\zusatzklammer {wobei die Operationen in $K$ ausgeführt werden} {} {.}}
\faktzusatz {}





}
{





Wir müssen zeigen, dass die konstante Nullfunktion und die konstante Einsfunktion auf $U$, das Negative einer algebraischen Funktion, und die Summe und das Produkt von zwei algebraischen Funktionen auf $U$ wieder algebraisch sind. Wir beschränken uns auf die Summe der algebraischen Funktionen $f_1$ und $f_2$. Es sei

\mavergleichskette
{\vergleichskette
{ P
}
{ \in }{ U
}
{  }{
}
{    }{
}
{   }{
}
}
{}{}{}
ein Punkt. Nach Voraussetzung gibt es Elemente

\mavergleichskette
{\vergleichskette
{ G_1,H_1,G_2,H_2
}
{ \in }{ R
}
{  }{
}
{    }{
}
{   }{
}
}
{}{}{}
mit

\mathdisp {f_1(Q) = \frac{G_1(Q)}{H_1(Q)} \text{ für alle }
 \mathdisplaybruch Q \in D(H_1) \subseteq U, P \in D(H_1)} { , }

und

\mathdisp {f_2(Q) = \frac{G_2(Q)}{H_2(Q)}  \text{ für alle }
 \mathdisplaybruch Q \in  D(H_2) \subseteq U, P \in D(H_2)} { . }

Es sei

\mavergleichskettedisp
{\vergleichskette
{ H
}
{ \defeq} { H_1 H_2
}
{ } {
}
{ } {
}
{ } {
}
}
{}{}{.}
Dann ist

\mavergleichskette
{\vergleichskette
{ P
}
{ \in }{ D(H)
}
{ = }{ D(H_1) \cap D(H_2)
}
{ \subseteq   }{ U
}
{   }{
}
}
{}{}{.}
Für einen beliebigen Punkt

\mavergleichskette
{\vergleichskette
{ Q
}
{ \in }{ D(H)
}
{  }{
}
{    }{
}
{   }{
}
}
{}{}{}
ist dann

\mavergleichskettealign
{\vergleichskettealign
{ (f_1+f_2)(Q)
}
{ =} { f_1(Q)+f_2(Q)
}
{ =} { \frac{G_1(Q)}{H_1(Q)} + \frac{G_2(Q)}{H_2(Q)}
}
{ =} { \frac{G_1(Q)H_2(Q) + G_2(Q) H_1(Q)}{H_1(Q) H_2(Q) }
}
{ =} { \frac{ (G_1H_2+G_2H_1)(Q) }{(H_1H_2)(Q)}
}
}
{}
{}{,}
was eine polynomiale Darstellung der Summenfunktion in der Zariski-offenen Umgebung
\mathl{D(H)}{} des Punktes $P$ ergibt.




}





\inputdefinition
{}
{





Es sei $K$ ein
\definitionsverweis {algebraisch abgeschlossener Körper}{}{,}
$R$ eine
$K$-\definitionsverweis {Algebra von endlichem Typ}{}{}
und sei

\mavergleichskette
{\vergleichskette
{ V
}
{ = }{ K\!-\!\operatorname{Spek}\, \left(  R  \right)
}
{  }{
}
{    }{
}
{   }{
}
}
{}{}{}
das
$K$-\definitionsverweis {Spektrum}{}{}
von $R$. Es sei

\mavergleichskette
{\vergleichskette
{ U
}
{ \subseteq }{ V
}
{  }{
}
{    }{
}
{   }{
}
}
{}{}{}
eine
\definitionsverweis {Zariski-offene Menge}{}{.}
Dann bezeichnet man mit

\mavergleichskettedisp
{\vergleichskette
{ \Gamma (U, {\mathcal O} )
}
{ =} { \left\{ f:U \longrightarrow K  \mid  f \text{ ist algebraisch}         \right\}
}
{ } {
}
{ } {
}
{ } {
}
}
{}{}{}
den \definitionswort {Ring der algebraischen Funktionen}{} auf $U$. Man bezeichnet ihn auch als \definitionswort {Strukturring zu}{} $U$ oder als \definitionswort {Schnittring}{} zu $U$.





}

Aufgrund von

Lemma 14.4

handelt es sich in der Tat um einen Ring. Das Symbol $\mathcal O$
\zusatzklammer {sprich \anfuehrung{O}{}} {} {}
bezeichnet die sogenannte \stichwort {Strukturgarbe} {.}





\inputfaktbeweis
{K-Spektrum/Ring der algebraischen Funktionen/Restriktion/Fakt}
{Lemma}
{}
{



\faktsituation {}
\faktvoraussetzung {Es sei $K$ ein
\definitionsverweis {algebraisch abgeschlossener Körper}{}{,}
$R$ eine $K$-
\definitionsverweis {Algebra von endlichem Typ}{}{}
und sei

\mavergleichskette
{\vergleichskette
{ V
}
{ = }{ K\!-\!\operatorname{Spek}\, \left(  R  \right)
}
{  }{
}
{    }{
}
{   }{
}
}
{}{}{}
das
$K$-\definitionsverweis {Spektrum}{}{}
von $R$. Es seien

\mavergleichskette
{\vergleichskette
{ U_1
}
{ \subseteq }{ U_2
}
{  }{
}
{    }{
}
{   }{
}
}
{}{}{}
offene Teilmengen von $V$.}
\faktfolgerung {Dann gibt es einen natürlichen
$K$-\definitionsverweis {Algebrahomomorphismus}{}{}
\maabbdisp {} { \Gamma (U_2 , {\mathcal O} ) } { \Gamma (U_1 , {\mathcal O} )
} {.}}
\faktzusatz {}
\faktzusatz {}





}
{





Die Funktion
\maabb {f} { U_2 } { K
} {}
liefert sofort durch Einschränkung eine auf $U_1$ definierte Funktion. Die lokal-algebraische Beschreibung, die für $f$ an jedem Punkt

\mavergleichskette
{\vergleichskette
{ P
}
{ \in }{ U_2
}
{  }{
}
{    }{
}
{   }{
}
}
{}{}{}
vorliegt, kann direkt auf der kleineren Teilmenge $U_1$ interpretiert werden.




}



Die im vorstehenden Lemma beschriebene Abbildung heißt \stichwort {Restriktionsabbildung} {} oder \stichwort {Einschränkungsabbildung } {.}





\inputfaktbeweis
{K-Spektrum/Ring der algebraischen Funktionen/U subseteq D(f)/Unabhängigkeit/Fakt}
{Lemma}
{}
{



\faktsituation {}
\faktvoraussetzung {Es sei $K$ ein
\definitionsverweis {algebraisch abgeschlossener Körper}{}{,}
$R$ eine
$K$-\definitionsverweis {Algebra von endlichem Typ}{}{}
und sei

\mavergleichskette
{\vergleichskette
{V
}
{ = }{ K\!-\!\operatorname{Spek}\, \left(  R  \right)
}
{  }{
}
{    }{
}
{   }{
}
}
{}{}{}
das
$K$-\definitionsverweis {Spektrum}{}{}
von $R$. Es sei

\mavergleichskette
{\vergleichskette
{ F
}
{ \in }{ R
}
{  }{
}
{    }{
}
{   }{
}
}
{}{}{}
und

\mavergleichskette
{\vergleichskette
{U
}
{ \subseteq }{ D(F)
}
{ \subseteq }{ V
}
{    }{
}
{   }{
}
}
{}{}{}
eine offene Menge.}
\faktfolgerung {Dann ist es egal, ob man
\mathl{\Gamma (U, {\mathcal O} )}{} mit Bezug auf $V$ oder mit Bezug auf

\mavergleichskette
{\vergleichskette
{ D(F)
}
{ = }{ K\!-\!\operatorname{Spek}\, \left(  R_F   \right)
}
{  }{
}
{    }{
}
{   }{
}
}
{}{}{}
definiert.}
\faktzusatz {}
\faktzusatz {}





}
{





Natürlich hängen die Funktionen auf $U$ nur von $U$ selbst ab, nicht von einem umgebenden Raum. Wir müssen zeigen, dass die lokal-algebraische Bedingung ebenfalls nur von $U$ abhängt. Es sei

\mavergleichskette
{\vergleichskette
{ P
}
{ \in }{ U
}
{  }{
}
{    }{
}
{   }{
}
}
{}{}{.}
Eine Beschreibung

\mathdisp {\varphi = \frac{G}{H} \text{ auf } D(H)
 \mathdisplaybruch \text{ mit } P \in D(H) \text{ und mit } G,H \in R} {  }

liefert sofort eine Beschreibung als Bruch auf
\mathl{D(HF)}{,} da man ja $H,G$ sofort in $R_F$ auffassen kann.

Es liege nun umgekehrt eine Bruchdarstellung

\mathdisp {\varphi = \frac{\tilde{G} }{\tilde{H} } \text{ auf } D(\tilde{H})
 \mathdisplaybruch \text{ mit } P \in D(\tilde{H}) \text{ und mit } \tilde{G},\tilde{H} \in R_F} {  }

vor. Es sei

\mavergleichskette
{\vergleichskette
{ \tilde{G}
}
{ = }{ G/F^r
}
{  }{
}
{    }{
}
{   }{
}
}
{}{}{}
und

\mavergleichskette
{\vergleichskette
{ \tilde{H}
}
{ = }{ H/F^s
}
{  }{
}
{    }{
}
{   }{
}
}
{}{}{.}
Dann gilt für jeden Punkt

\mavergleichskette
{\vergleichskette
{ Q
}
{ \in }{ D(HF)
}
{  }{
}
{    }{
}
{   }{
}
}
{}{}{}
die Gleichheit

\mavergleichskettedisp
{\vergleichskette
{ \varphi(Q)
}
{ =} { \frac{\tilde{G}(Q)}{\tilde{H}(Q)}
}
{ =} { \frac{G(Q)/F^r(Q)}{ H(Q)/F^s(Q)}
}
{ =} { \frac{G(Q)F^s(Q)}{ H(Q)F^r(Q)}
}
{ } {
}
}
{}{}{.}
Dabei haben wir im letzten Schritt mit
\mathl{F^{r+s}}{} erweitert. In der letzten Darstellung sind Zähler und Nenner aus $R$, und es ist

\mavergleichskette
{\vergleichskette
{ HF^r(P)
}
{ \neq }{ 0
}
{  }{
}
{    }{
}
{   }{
}
}
{}{}{,}
also ist
\mathl{D(HF^r)}{} eine offene Umgebung von $P$.




}






\inputfaktbeweis
{K-Spektrum/Algebraisch abgeschlossener Körper/Verschiedene rationale Darstellungen einer aIgebraischen Funktion/Beziehung im Koordinatenring/Fakt}
{Lemma}
{}
{



\faktsituation {}
\faktvoraussetzung {Es sei $K$ ein
\definitionsverweis {algebraisch abgeschlossener Körper}{}{,}
$R$ eine $K$-
\definitionsverweis {Algebra von endlichem Typ}{}{}
und sei

\mavergleichskette
{\vergleichskette
{ V
}
{ = }{ K\!-\!\operatorname{Spek}\, \left(  R  \right)
}
{  }{
}
{    }{
}
{   }{
}
}
{}{}{}
das
$K$-\definitionsverweis {Spektrum}{}{}
von $R$. Es sei

\mavergleichskette
{\vergleichskette
{ U
}
{ \subseteq }{ V
}
{  }{
}
{    }{
}
{   }{
}
}
{}{}{}
eine Zariski-offene Menge,

\mavergleichskette
{\vergleichskette
{ P
}
{ \in }{ U
}
{  }{
}
{    }{
}
{   }{
}
}
{}{}{}
ein Punkt und es sei
\maabb {f} { U } { K
} {}
eine
\definitionsverweis {algebraische Funktion}{}{,}
für die es die beiden rationalen Darstellungen

\mathdisp {\frac{G_1 }{H_1 } \text{ und } \frac{G_2 }{H_2 }} {  }

gebe mit

\mavergleichskette
{\vergleichskette
{ G_1,H_1,G_2,H_2
}
{ \in }{ R
}
{  }{
}
{    }{
}
{   }{
}
}
{}{}{}
und mit

\mavergleichskette
{\vergleichskette
{ P
}
{ \in }{ D(H_1), D(H_2)
}
{ \subseteq }{ U
}
{    }{
}
{   }{
}
}
{}{}{.}}
\faktfolgerung {Dann gibt es ein

\mavergleichskette
{\vergleichskette
{ r
}
{ \in }{ \N
}
{  }{
}
{    }{
}
{   }{
}
}
{}{}{}
mit

\mathdisp {H_1^rH_2^r(G_1H_2-G_2H_1)^r = 0 \text{ in } R} { . }
}
\faktzusatz {Ist $R$
\definitionsverweis {reduziert}{}{,}
so gilt sogar

\mavergleichskettedisp
{\vergleichskette
{ H_1H_2(G_1H_2-G_2H_1)
}
{ =} { 0
}
{ } {
}
{ } {
}
{ } {
}
}
{}{}{.}}
\faktzusatz {}





}
{





Wir betrachten das Element

\mavergleichskette
{\vergleichskette
{ F
}
{ = }{ H_1H_2 \left(  G_1H_2-G_2H_1  \right)
}
{  }{
}
{    }{
}
{   }{
}
}
{}{}{}
auf $V$ und behaupten, dass dies die Nullfunktion induziert. Es sei

\mavergleichskette
{\vergleichskette
{ Q
}
{ \in }{ V
}
{  }{
}
{    }{
}
{   }{
}
}
{}{}{.}
Bei

\mavergleichskette
{\vergleichskette
{ H_1(Q)
}
{ = }{ 0
}
{  }{
}
{    }{
}
{   }{
}
}
{}{}{}
oder

\mavergleichskette
{\vergleichskette
{ H_2(Q)
}
{ = }{ 0
}
{  }{
}
{    }{
}
{   }{
}
}
{}{}{}
ist

\mavergleichskette
{\vergleichskette
{ F(Q)
}
{ = }{ 0
}
{  }{
}
{    }{
}
{   }{
}
}
{}{}{,}
sei also

\mavergleichskette
{\vergleichskette
{ H_1(Q), H_2(Q)
}
{ \neq }{ 0
}
{  }{
}
{    }{
}
{   }{
}
}
{}{}{}
vorausgesetzt. Dann ist

\mavergleichskette
{\vergleichskette
{ Q
}
{ \in }{ D(H_1) \cap D(H_2)
}
{  }{
}
{    }{
}
{   }{
}
}
{}{}{,}
und dort gelten die beiden rationalen Darstellungen für $f$, nämlich

\mavergleichskettedisp
{\vergleichskette
{ \frac{G_1(Q)}{H_1(Q)}
}
{ =} { f(Q)
}
{ =} { \frac{G_2(Q)}{H_2(Q)}
}
{ } {
}
{ } {
}
}
{}{}{.}
Daraus folgt

\mavergleichskette
{\vergleichskette
{ G_1(Q)H_2(Q)
}
{ = }{ G_2(Q)H_1(Q)
}
{  }{
}
{    }{
}
{   }{
}
}
{}{}{}
und somit ist die Differenz $0$. Insgesamt ist also $F$ die Nullfunktion auf $V$ und daher gibt es nach
dem
Hilbertschen Nullstellensatz

ein $r$ mit

\mavergleichskette
{\vergleichskette
{ F^r
}
{ = }{ 0
}
{  }{
}
{    }{
}
{   }{
}
}
{}{}{}
in $R$.




}




\bild{ \begin{center}
\includegraphics[width=5.5cm]{ \bildeinlesungpng {Monkey_Saddle_Surface_(Shaded)} {png} }
\end{center}
\bildtext {Der Graph einer globalen Funktion auf ${\mathbb A}^{2}_{K}$.} }
\bildlizenz { Monkey Saddle Surface (Shaded).png } {} {Inductiveload} {Commons} {PD} {}







\inputfaktbeweis
{K-Spektrum/Algebraisch abgeschlossener Körper/Algebraische (reguläre) Funktion auf offener Menge/Globaler Schnittring ist Koordinatenring/Fakt}
{Satz}
{}
{



\faktsituation {}
\faktvoraussetzung {Es sei $K$ ein
\definitionsverweis {algebraisch abgeschlossener Körper}{}{,}
$R$ eine
\definitionsverweis {reduzierte}{}{}
$K$-\definitionsverweis {Algebra von endlichem Typ}{}{}
und sei

\mavergleichskette
{\vergleichskette
{ V
}
{ = }{ K\!-\!\operatorname{Spek}\, \left(  R  \right)
}
{  }{
}
{    }{
}
{   }{
}
}
{}{}{}
das
$K$-\definitionsverweis {Spektrum}{}{}
von $R$.}
\faktfolgerung {Dann ist

\mavergleichskettedisp
{\vergleichskette
{ \Gamma (V, {\mathcal O} )
}
{ =} { R
}
{ } {
}
{ } {
}
{ } {
}
}
{}{}{.}}
\faktzusatz {}
\faktzusatz {}





}
{





Ein Element

\mavergleichskette
{\vergleichskette
{ F
}
{ \in }{ R
}
{  }{
}
{    }{
}
{   }{
}
}
{}{}{}
liefert direkt eine algebraische Funktion auf ganz $V$, was einen
$K$-\definitionsverweis {Algebrahomomorphismus}{}{}
\maabbdisp {} { R } { \Gamma (V, {\mathcal O} )
} {}
ergibt. Wenn dabei $F$ an jedem Punkt die Nullfunktion induziert, so ist nach

Satz 11.1

und wegen der Reduziertheit auch

\mavergleichskette
{\vergleichskette
{ F
}
{ = }{ 0
}
{  }{
}
{    }{
}
{   }{
}
}
{}{}{.}
D.h. die Abbildung ist injektiv.

Es sei nun
\maabb {f} { V } { K
} {}
ein algebraische Funktion. Dann gibt es zu jedem Punkt

\mavergleichskette
{\vergleichskette
{ P
}
{ \in }{ V
}
{  }{
}
{    }{
}
{   }{
}
}
{}{}{}
Elemente

\mavergleichskette
{\vergleichskette
{ G_P, H_P
}
{ \in }{ R
}
{  }{
}
{    }{
}
{   }{
}
}
{}{}{}
mit

\mavergleichskette
{\vergleichskette
{ P
}
{ \in }{ D(H_P)
}
{  }{
}
{    }{
}
{   }{
}
}
{}{}{}
und mit

\mavergleichskette
{\vergleichskette
{ f
}
{ = }{ { \frac{  G_P  }{  H_P  } }
}
{  }{
}
{    }{
}
{   }{
}
}
{}{}{}
auf
\mathl{D(H_P)}{.} Die $D(H_P)$ bilden eine offene Überdeckung von $V$ und das bedeutet nach

Korollar 11.12
,
dass die $H_P$ in $R$ das
\definitionsverweis {Einheitsideal}{}{}
erzeugen. Dann gibt es aber auch eine endliche Auswahl davon, die das Einheitsideal erzeugen, sagen wir

\mathbed {H_i = H_{P_i }} {}
 {i= 1 , \ldots ,  m} {}
 {} {} {} {.}
Dann wiederum überdecken diese

\mathbed {D(H_i)} {}
 {i= 1 , \ldots ,  m} {}
 {} {} {} {,}
ganz $V$.

Auf den Durchschnitten

\mavergleichskette
{\vergleichskette
{ D(H_iH_j)
}
{ = }{ D(H_i) \cap D(H_j)
}
{  }{
}
{    }{
}
{   }{
}
}
{}{}{}
haben wir die Identitäten

\mathdisp {f(Q) = \frac{G_i(Q)}{H_i(Q)} = \frac{G_j(Q)}{H_j(Q)} \text{ für alle } Q \in D(H_iH_j)} { . }

Daraus folgt nach

Lemma 14.8

und der Reduziertheit, dass

\mavergleichskettedisp
{\vergleichskette
{ H_iH_j G_iH_j
}
{ =} { H_iH_j G_j H_i
}
{ } {
}
{ } {
}
{ } {
}
}
{}{}{}
in $R$ gilt. Wir ersetzen $H_i$ durch $H_i^2$ und $G_i$ durch
\mathl{G_i H_i}{.} Dann ist nach wie vor $G_i/H_i$ eine lokale Beschreibung für $f$, und die letzte Bedingung vereinfacht sich zu

\mavergleichskettedisp
{\vergleichskette
{ H_i G_j
}
{ =} { H_j G_i
}
{ } {
}
{ } {
}
{ } {
}
}
{}{}{.}

Da die $H_i$ das Einheitsideal erzeugen, gibt es Elemente

\mavergleichskette
{\vergleichskette
{ A_i
}
{ \in }{ R
}
{  }{
}
{    }{
}
{   }{
}
}
{}{}{}
mit

\mavergleichskettedisp
{\vergleichskette
{ \sum_{i = 1}^m A_i H_i
}
{ =} { 1
}
{ } {
}
{ } {
}
{ } {
}
}
{}{}{}
in $R$. Wir behaupten, dass das Element

\mavergleichskettedisp
{\vergleichskette
{ F
}
{ =} { \sum_{i = 1}^m A_i G_i
}
{ } {
}
{ } {
}
{ } {
}
}
{}{}{}
auf ganz $V$ die Funktion $f$ induziert. Dazu sei

\mavergleichskette
{\vergleichskette
{ Q
}
{ \in }{ V
}
{  }{
}
{    }{
}
{   }{
}
}
{}{}{}
ein beliebiger Punkt, und zwar sei ohne Einschränkung

\mavergleichskette
{\vergleichskette
{ Q
}
{ \in }{ D(H_1)
}
{  }{
}
{    }{
}
{   }{
}
}
{}{}{.}
Dann ist

\mavergleichskettealign
{\vergleichskettealign
{ f(Q)
}
{ =} { \frac{G_1(Q)}{H_1(Q)}
}
{ =} { \frac{G_1(Q)}{H_1(Q)} \left(  \sum_{i = 1}^m A_i H_i  \right) (Q)
}
{ =} { \sum_{i = 1}^m A_i(Q) \frac{ G_1(Q) H_i(Q) }{H_1(Q)}
}
{ =} { \sum_{i = 1}^m A_i(Q) G_i(Q)
}
}
{
\vergleichskettefortsetzungalign
{ =} { F(Q)
}
{ } {}
{ } {}
{ } {}
}
{}{.}




}






\inputfaktbeweis
{K-Spektrum/Algebraisch abgeschlossener Körper/Algebraische (reguläre) Funktion auf D(f)/Ist R_f/Fakt}
{Korollar}
{}
{



\faktsituation {}
\faktvoraussetzung {Es sei $K$ ein
\definitionsverweis {algebraisch abgeschlossener Körper}{}{,}
$R$ eine
\definitionsverweis {reduzierte}{}{}
$K$-
\definitionsverweis {Algebra von endlichem Typ}{}{}
und sei

\mavergleichskette
{\vergleichskette
{ V
}
{ = }{ K\!-\!\operatorname{Spek}\, \left(  R  \right)
}
{  }{
}
{    }{
}
{   }{
}
}
{}{}{}
das
$K$-\definitionsverweis {Spektrum}{}{}
von $R$. Es sei

\mavergleichskette
{\vergleichskette
{ F
}
{ \in }{ R
}
{  }{
}
{    }{
}
{   }{
}
}
{}{}{}
mit zugehöriger offener Menge

\mavergleichskette
{\vergleichskette
{ D(F)
}
{ \subseteq }{ V
}
{  }{
}
{    }{
}
{   }{
}
}
{}{}{.}}
\faktfolgerung {Dann ist

\mavergleichskettedisp
{\vergleichskette
{ \Gamma (D(F), {\mathcal O} )
}
{ =} { R_F
}
{ } {
}
{ } {
}
{ } {
}
}
{}{}{.}}
\faktzusatz {}
\faktzusatz {}





}
{





Dies folgt direkt aus

Lemma 14.7

und

Satz 14.9
.




}


\inputbemerkung
{}
{





Eine Variante des 14. Hilbertschen Problems ist die Frage, ob zu jeder offenen Menge $U$ der Ring der
\definitionsverweis {algebraischen Funktionen}{}{}
$\Gamma ( U , {\mathcal O} )$
\definitionsverweis {endlich erzeugt}{}{}
ist. Für offene Mengen der Form

\mavergleichskette
{\vergleichskette
{U
}
{ = }{ D(f)
}
{  }{
}
{    }{
}
{   }{
}
}
{}{}{}
ist dies richtig, ebenfalls, wenn $R$ regulär oder faktoriell ist, und auch in kleinen Dimensionen. Im Allgemeinen ist es aber nicht richtig.





}
